\documentclass[11pt,a4paper]{amsart}

\usepackage{amsmath,amssymb,amsthm}
\usepackage[margin=1in]{geometry}
\usepackage{hyperref}
\usepackage{enumerate}

% Theorem environments
\newtheorem{theorem}{Theorem}[section]
\newtheorem{lemma}[theorem]{Lemma}
\newtheorem{proposition}[theorem]{Proposition}
\newtheorem{corollary}[theorem]{Corollary}
\newtheorem{conjecture}[theorem]{Conjecture}

\theoremstyle{definition}
\newtheorem{definition}[theorem]{Definition}
\newtheorem{example}[theorem]{Example}
\newtheorem{remark}[theorem]{Remark}

% Macros
\newcommand{\D}{\mathbb{D}}
\newcommand{\T}{\mathbb{T}}
\newcommand{\C}{\mathbb{C}}
\newcommand{\R}{\mathbb{R}}
\newcommand{\Pn}{\mathcal{P}_n}
\newcommand{\abs}[1]{\lvert #1 \rvert}
\newcommand{\norm}[1]{\lVert #1 \rVert}
\newcommand{\diam}{\operatorname{diam}}
\newcommand{\inrad}{\operatorname{inrad}}
\renewcommand{\Cap}{\operatorname{Cap}}

\title[Inradius of polynomial lemniscates]{On the inradius of polynomial lemniscates: \\
the non-degenerate case and beyond}

\author{Tran Duc Huy}

\address{SoznaLab}

\date{\today}

\begin{document}

\begin{abstract}
We prove that for any monic polynomial $P$ of degree $n$ with all roots in
the closed unit disc $\overline{\D}$, if every critical value of $P$ has
modulus at least $1$, then the lemniscate $\{z \in \C : \abs{P(z)} < 1\}$
has inradius at least $1/(4n)$.  This confirms the Erd\H{o}s--Herzog--Piranian
conjecture ($\rho_n \gg 1/n$) for the natural class of \emph{non-degenerate}
polynomials, which includes the conjectured extremizer $P(z) = z^n - 1$.
The proof combines Koebe's quarter theorem with a Fekete energy
selection argument and is independent of the recent area-based approach of
Krishnapur--Lundberg--Ramachandran (2025), who applied
a Nazarov--Polterovich--Sodin sublevel set bound to establish
$\rho_n \geq c/(n\sqrt{\log n})$ for all polynomials.
We further analyze the squarefree degenerate case via Blaschke product
factorizations, identify an obstruction to extending the conformal
method directly to merged components, and show that the full conjecture
would follow if an inradius minimizer were non-degenerate.
Preliminary numerics are consistent with this scenario.
\end{abstract}

\maketitle

%% ====================================================================
\section{Introduction}
%% ====================================================================

For a monic polynomial $P(z) = \prod_{i=1}^n (z - z_i)$ with coefficients in $\C$,
the \emph{lemniscate} at level $t > 0$ is
\[
  \Lambda_P(t) \;=\; \bigl\{z \in \C : \abs{P(z)} < t\bigr\}.
\]
When all roots satisfy $\abs{z_i} \leq 1$, we write
$P \in \Pn(\overline{\D})$.
The \emph{inradius} of a bounded open set $\Omega \subset \C$ is
\[
  \inrad(\Omega) \;=\; \sup\bigl\{r > 0 : B(z,r) \subset \Omega
  \text{ for some } z \in \C\bigr\}.
\]

In their foundational 1958 paper \cite{EHP58}, Erd\H{o}s, Herzog, and Piranian
posed the problem of determining the minimal inradius
\[
  \rho_n \;=\; \inf\bigl\{\inrad(\Lambda_P(1)) : P \in \Pn(\overline{\D})\bigr\}.
\]
The polynomial $P(z) = z^n - 1$ has $n$ roots at the $n$-th roots of unity
and yields $\inrad(\Lambda_{z^n-1}) \sim \pi/(2n)$, suggesting the conjecture:

\begin{conjecture}[Erd\H{o}s--Herzog--Piranian]
$\rho_n \asymp 1/n$.
\end{conjecture}

Pommerenke \cite{Pom61} established the first lower bound
$\rho_n \geq 1/(2en^2)$ in 1961, which remained the state of the art for
64 years.  In a remarkable 2025 breakthrough, Krishnapur, Lundberg, and
Ramachandran \cite{KLR25} proved
\begin{equation}\label{eq:KLR}
  \rho_n \;\geq\; \frac{c}{n\sqrt{\log n}},
\end{equation}
confirming the conjecture up to a factor of $\sqrt{\log n}$.
Their proof applies a Nazarov--Polterovich--Sodin lower bound on the area of
harmonic sublevel sets with an isoperimetric inequality relating area,
perimeter, and inradius of lemniscate components.  The $\sqrt{\log n}$ loss
arises from the square root in the area-to-inradius conversion.
(Separately, Tao~\cite{Tao25} recently resolved the related
\emph{length} conjecture from the same 1958 paper, establishing that the
maximal arclength of $\{\abs{P(z)} = 1\}$ is $\Theta(n)$ for
sufficiently large $n$.)

In this note, we take a different approach.  We identify a natural class of
polynomials---those whose critical values all lie outside the open unit
disc---for which the sharp $1/n$ bound holds.  Our proof bypasses the area
estimates entirely, using instead the univalent branch structure of $P$
together with a classical energy bound.  We then analyze the degenerate
case (where some critical values lie inside $\D$) via Blaschke product
factorizations, identify a structural obstruction that prevents na\"{\i}ve
extensions of the conformal method, and show that the full conjecture
would follow if some inradius minimizer were non-degenerate.

%% ====================================================================
\section{Definitions and statement of results}
%% ====================================================================

\begin{definition}
A monic polynomial $P \in \Pn(\overline{\D})$ is called \emph{non-degenerate}
if every critical value of $P$ has modulus at least $1$:
\[
  P'(\zeta) = 0 \quad\Longrightarrow\quad \abs{P(\zeta)} \geq 1.
\]
Let $\Pn^{\mathrm{nd}}(\overline{\D})$ denote the set of non-degenerate
polynomials in $\Pn(\overline{\D})$.
\end{definition}

The non-degeneracy condition means that the polynomial $P$, viewed as a
branched covering $P \colon \C \to \C$, has no branch points over the open
unit disc $\D$.  Equivalently, every branch of $P^{-1}$ extends
univalently from $0$ to the boundary $\abs{w} = 1$.

\begin{example}\label{ex:roots-of-unity}
The polynomial $P(z) = z^n - 1$ has a unique critical point at $\zeta = 0$,
with critical value $P(0) = -1$ and $\abs{P(0)} = 1$.  Thus
$z^n - 1 \in \Pn^{\mathrm{nd}}(\overline{\D})$.
\end{example}

\begin{example}\label{ex:simple-roots}
More generally, if $P$ has $n$ distinct roots and all roots lie on $\T$,
then $P \in \Pn^{\mathrm{nd}}(\overline{\D})$ whenever
$\abs{P(\zeta)} \geq 1$ for each of the $n-1$ critical points.
Numerical evidence suggests this holds for a wide class of polynomials
with roots on or near the unit circle.
\end{example}

Our main result is:

\begin{theorem}[Inradius bound for non-degenerate lemniscates]
\label{thm:main}
Let $P \in \Pn^{\mathrm{nd}}(\overline{\D})$.  Then
\[
  \inrad(\Lambda_P(1)) \;\geq\; \frac{1}{4n}.
\]
\end{theorem}

Since $z^n - 1$ is non-degenerate with $\inrad(\Lambda_{z^n-1}) \sim \pi/(2n)$
(Proposition~\ref{prop:extremal}),
the bound has the sharp order $1/n$.

\begin{remark}
The constant $1/4$ arises from the Koebe quarter theorem and is not optimal.
For the extremal polynomial $z^n - 1$ the true constant is $\pi/2 \approx 1.571$.
The Koebe bound is sharp for the full class of univalent functions, but
the inverse branches $\psi_j$ satisfy additional constraints
(e.g., $\psi_j(\D)$ is a petal of a polynomial lemniscate).
Whether refined distortion theorems or Bieberbach-type
coefficient bounds for $\psi_j$ can improve the constant for
$\Pn^{\mathrm{nd}}(\overline{\D})$ remains an interesting question.
\end{remark}

As an immediate consequence, we recover the EHP conjecture for the
non-degenerate class:

\begin{corollary}\label{cor:nd-class}
$\inf\bigl\{\inrad(\Lambda_P(1)) : P \in \Pn^{\mathrm{nd}}(\overline{\D})\bigr\}
\asymp 1/n$.
\end{corollary}

%% ====================================================================
\section{Proof of Theorem~\ref{thm:main}}
%% ====================================================================

The proof has three components: univalent branch extraction,
the Koebe quarter theorem, and a Fekete energy argument for root selection.

%% --------------------------------------------------
\subsection{Univalent inverse branches}
%% --------------------------------------------------

\begin{lemma}\label{lem:branch}
Let $P \in \Pn^{\mathrm{nd}}(\overline{\D})$ with distinct roots
$z_1, \ldots, z_n$.  For each $j = 1, \ldots, n$, there exists a
holomorphic function $\psi_j \colon \D \to \C$ such that:
\begin{enumerate}[(i)]
  \item $P(\psi_j(w)) = w$ for all $w \in \D$;
  \item $\psi_j(0) = z_j$;
  \item $\psi_j$ is univalent on $\D$;
  \item $\psi_j(\D) \subset \Lambda_P(1)$.
\end{enumerate}
\end{lemma}

\begin{proof}
Since $P(z_j) = 0$ and $P'(z_j) \neq 0$ (non-degeneracy implies simple
roots by Remark~\ref{rem:repeated}), the inverse function theorem provides a local
holomorphic branch $\psi_j$ near $w = 0$ with $\psi_j(0) = z_j$.

The branch $\psi_j$ can be continued analytically along any path in $\D$
that avoids the critical values of $P$.  By the non-degeneracy hypothesis,
all critical values have modulus $\geq 1$ and therefore lie outside $\D$.
By the monodromy theorem, $\psi_j$ extends to a single-valued holomorphic
function on all of $\D$.

Univalence follows because $\psi_j(w_1) = \psi_j(w_2)$ would imply
$w_1 = P(\psi_j(w_1)) = P(\psi_j(w_2)) = w_2$.

Finally, $w \in \D$ implies $\abs{P(\psi_j(w))} = \abs{w} < 1$, so
$\psi_j(w) \in \Lambda_P(1)$.
\end{proof}

\begin{remark}\label{rem:repeated}
If $P$ has a repeated root at $z_j$ of multiplicity $m \geq 2$, then
$z_j$ is itself a critical point with critical value $P(z_j) = 0$.
The non-degeneracy condition then requires $\abs{P(z_j)} \geq 1$,
but $P(z_j) = 0$, a contradiction.  Therefore, all roots of a
non-degenerate polynomial are simple.
\end{remark}

%% --------------------------------------------------
\subsection{The Koebe quarter theorem}
%% --------------------------------------------------

We recall the classical result (see e.g.\ \cite[Chapter~2]{Duren83}):

\begin{theorem}[Koebe, 1907]\label{thm:koebe}
If $f \colon \D \to \C$ is univalent with $f(0) = a$ and $f'(0) = b$, then
\[
  f(\D) \;\supset\; B\!\left(a,\, \tfrac{\abs{b}}{4}\right).
\]
\end{theorem}

Applying Theorem~\ref{thm:koebe} to $\psi_j$ from Lemma~\ref{lem:branch}:

\begin{corollary}\label{cor:petal-disc}
For each root $z_j$,
\[
  \Lambda_P(1) \;\supset\; \psi_j(\D) \;\supset\;
  B\!\left(z_j,\, \frac{1}{4\abs{P'(z_j)}}\right).
\]
\end{corollary}

\begin{proof}
We have $\psi_j(0) = z_j$ and $\psi_j'(0) = 1/P'(z_j)$, the latter
following from differentiating $P(\psi_j(w)) = w$ at $w = 0$.
\end{proof}

%% --------------------------------------------------
\subsection{Root selection via Fekete energy}
%% --------------------------------------------------

To complete the proof, we must show that at least one root has
$\abs{P'(z_j)} \leq n$.

\begin{lemma}[Fekete energy bound]\label{lem:fekete}
Let $z_1, \ldots, z_n \in \overline{\D}$ be distinct.  Then
\[
  \min_{1 \leq j \leq n} \abs{P'(z_j)} \;\leq\; n.
\]
\end{lemma}

\begin{proof}
Since $P(z) = \prod_{i=1}^n (z - z_i)$, we have
$P'(z_j) = \prod_{i \neq j} (z_j - z_i)$, and therefore
\[
  \prod_{j=1}^n \abs{P'(z_j)}
  \;=\; \prod_{j=1}^n \prod_{i \neq j} \abs{z_j - z_i}
  \;=\; \Bigl(\prod_{1 \leq i < j \leq n} \abs{z_i - z_j}\Bigr)^2.
\]
Consider the Vandermonde matrix $V = (z_i^{k-1})_{1 \leq i,k \leq n}$,
so that $\abs{\det V} = \prod_{1 \leq i < j \leq n}\abs{z_j - z_i}$.
By Hadamard's inequality, $\abs{\det V} \leq \prod_{i=1}^n \norm{r_i}_2$,
where $r_i = (1, z_i, \ldots, z_i^{n-1})$.  Since $\abs{z_i} \leq 1$,
\[
  \norm{r_i}_2^2 \;=\; \sum_{k=0}^{n-1}\abs{z_i}^{2k} \;\leq\; n,
\]
so $\abs{\det V} \leq n^{n/2}$.  Therefore
\[
  \prod_{j=1}^n \abs{P'(z_j)} \;\leq\; n^n.
\]
The geometric mean of the $\abs{P'(z_j)}$ is at most $n$,
so $\min_j \abs{P'(z_j)} \leq n$.

Equality is attained for the $n$-th roots of unity, where
$\abs{P'(\omega_k)} = n$ for every $k$, confirming sharpness.
\end{proof}

%% --------------------------------------------------
\subsection{Completion of the proof}
%% --------------------------------------------------

\begin{proof}[Proof of Theorem~\ref{thm:main}]
Let $P \in \Pn^{\mathrm{nd}}(\overline{\D})$.  By Remark~\ref{rem:repeated},
all roots are simple.  Choose $j^*$ attaining $\min_j \abs{P'(z_j)}$.
By Lemma~\ref{lem:fekete}, $\abs{P'(z_{j^*})} \leq n$.
By Corollary~\ref{cor:petal-disc},
\[
  \Lambda_P(1) \;\supset\;
  B\!\left(z_{j^*},\, \frac{1}{4\abs{P'(z_{j^*})}}\right)
  \;\supset\;
  B\!\left(z_{j^*},\, \frac{1}{4n}\right).
\]
Therefore $\inrad(\Lambda_P(1)) \geq 1/(4n)$.
\end{proof}

%% ====================================================================
\section{Verification on the extremal polynomial}
%% ====================================================================

\begin{proposition}\label{prop:extremal}
For $P_n(z) = z^n - 1$, one has
\[
  \inrad(\Lambda_{P_n}(1)) \;\sim\; \frac{\pi}{2n}
  \qquad\text{as } n \to \infty.
\]
More precisely,
\[
  \liminf_{n \to \infty}\, n \cdot \inrad(\Lambda_{P_n}(1)) \;\geq\; \frac{\pi}{2}
  \qquad\text{and}\qquad
  \inrad(\Lambda_{P_n}(1)) \;\leq\; \frac{\pi}{2n} + O\!\left(\frac{1}{n^2}\right).
\]
In particular, the bound of Theorem~\ref{thm:main} has the sharp order $1/n$.
\end{proposition}

\begin{proof}
By rotational symmetry, $\Lambda_{P_n}(1)$ is the union of $n$ congruent
components, so it suffices to analyze the component $U_n$ containing $1$.
On $U_n$ we may use the principal logarithm and write $z = e^{\xi/n}$.
Then
\[
  z \in U_n \iff \abs{e^\xi - 1} < 1 \text{ and } \abs{\mathrm{Im}\,\xi} < \pi/2.
\]
Hence $U_n = \exp(\Omega/n)$, where
\[
  \Omega \;:=\; \{\xi = x + iy : x < \log(2\cos y),\; \abs{y} < \pi/2\}.
\]

We first show that $\inrad(\Omega) = \pi/2$.
Since $\Omega \subset \{\abs{y} < \pi/2\}$, every disc contained in $\Omega$
has radius at most $\pi/2$, so $\inrad(\Omega) \leq \pi/2$.
Conversely, let $0 < r < \pi/2$.  Choose $M > r - \log(2\cos r)$.
If $\abs{\xi + M} < r$, then $\abs{\mathrm{Im}\,\xi} < r$ and
\[
  \mathrm{Re}\,\xi \;\leq\; -M + r \;<\; \log(2\cos r)
  \;\leq\; \log(2\cos(\mathrm{Im}\,\xi)),
\]
so $B(-M, r) \subset \Omega$.  Therefore $\inrad(\Omega) \geq r$ for every
$r < \pi/2$, proving $\inrad(\Omega) = \pi/2$.

\emph{Lower bound on $\inrad(U_n)$.}
Fix $0 < r < \pi/2$ and choose $M$ as above so that
$B(-M, r) \subset \Omega$.  Set $a_n = e^{-M/n}$.
For $\abs{h} \leq r$,
\[
  e^{(-M+h)/n} \;=\; a_n\!\left(1 + \frac{h}{n} + R_n(h)\right),
\]
where $\abs{R_n(h)} \leq C_r/n^2$ for some constant $C_r$ depending
only on $r$.  Fix $\varepsilon \in (0, r)$.  If
$\abs{w - a_n} < a_n(r - \varepsilon)/n$, set
\[
  F(h) \;=\; a_n \frac{h}{n} - (w - a_n), \qquad
  G_n(h) \;=\; a_n R_n(h).
\]
On $\abs{h} = r$ we have $\abs{F(h)} \geq a_n \varepsilon/n$ and
$\abs{G_n(h)} \leq a_n C_r/n^2$.  Hence for all sufficiently large $n$,
$\abs{G_n(h)} < \abs{F(h)}$ on $\abs{h} = r$.
By Rouch\'{e}'s theorem, $F + G_n$ has a zero in $\abs{h} < r$, i.e.\
there exists $\abs{h} < r$ such that $e^{(-M+h)/n} = w$.  Therefore
\[
  B\!\left(a_n,\, a_n\frac{r - \varepsilon}{n}\right) \;\subset\; U_n
\]
for all large $n$.  Since $a_n = 1 + O(1/n)$,
$\liminf_{n \to \infty} n \cdot \inrad(U_n) \geq r - \varepsilon$.
Letting $\varepsilon \downarrow 0$ and then $r \uparrow \pi/2$,
\[
  \liminf_{n \to \infty} n \cdot \inrad(U_n) \;\geq\; \frac{\pi}{2}.
\]

\emph{Upper bound on $\inrad(U_n)$.}
If $z \in U_n$, then $\abs{z}^n < 2$ and $\abs{\arg z} < \pi/(2n)$.
Thus $U_n$ is contained in the sector
$\{z : \abs{\arg z} < \pi/(2n)\}$ and in $\{\abs{z} < 2^{1/n}\}$.
If $B(z_0, \rho) \subset U_n$, then $\rho$ is at most the distance from
$z_0$ to the boundary of the sector, hence
\[
  \rho \;\leq\; \abs{z_0}\sin\!\left(\frac{\pi}{2n}\right)
  \;\leq\; 2^{1/n}\sin\!\left(\frac{\pi}{2n}\right)
  \;=\; \frac{\pi}{2n} + O\!\left(\frac{1}{n^2}\right).
\]

Combining the bounds gives
$\inrad(U_n) \leq \pi/(2n) + O(1/n^2)$ and
$\liminf_{n \to \infty} n \cdot \inrad(U_n) \geq \pi/2$,
which together yield $\inrad(\Lambda_{P_n}(1)) \sim \pi/(2n)$.
\end{proof}

%% ====================================================================
\section{The non-degenerate class}
%% ====================================================================

The non-degeneracy condition $\abs{P(\zeta)} \geq 1$ at all critical points
has a clean potential-theoretic interpretation.  The logarithmic potential
$u(z) = \log\abs{P(z)}$ satisfies $u(z_j) = -\infty$ at the roots
and $u(\zeta) \geq 0$ at the critical points.
Away from the roots, the zeros of $P'$ are exactly the critical points
of $u$; since $u$ is harmonic there, simple critical points are of
saddle type.
Non-degeneracy means that $u$ is non-negative at all its finite
critical points.

\begin{lemma}\label{lem:petals}
$P$ is non-degenerate if and only if $\Lambda_P(1)$ has
exactly $n$ connected components, each containing a single simple root.
\end{lemma}

\begin{proof}
Every connected component $U$ of $\Lambda_P(1)$ is simply connected.
To see this, let $\gamma \subset U$ be a simple closed curve bounding
a domain $D_\gamma$.  Since $\abs{P} < 1$ on $\gamma$ and $\abs{P}$
is subharmonic (as the modulus of a holomorphic function), the maximum
principle for subharmonic functions gives $\abs{P} < 1$ on $D_\gamma$.
Thus $D_\gamma \subset \Lambda_P(1)$, and since $D_\gamma$ is connected
and meets $U$ (along $\gamma$), we have $D_\gamma \subset U$.
Therefore $U$ has no bounded complementary components.

The restriction $P|_U \colon U \to \D$ is a proper holomorphic map
whose degree equals the number of roots in $U$ counted with multiplicity.

$(\Rightarrow)$  If $P$ is non-degenerate, $P$ has no critical values in
$\D$, so $P|_U$ is an unbranched covering of $\D$.  Since $\D$ is simply
connected, this covering has degree $1$: each component contains exactly
one simple root, and there are $n$ components.

$(\Leftarrow)$  If $\Lambda_P(1)$ has $n$ components each with one root,
then $P|_U$ has degree $1$ on each component, hence no critical points
in $\Lambda_P(1)$.  Every critical value therefore has modulus $\geq 1$.
\end{proof}

This structural property---that the lemniscate has $n$ separated petals---is
essential to our proof and is the geometric content of non-degeneracy.

%% ====================================================================
\section{Analysis of the degenerate case}
\label{sec:degenerate}
%% ====================================================================

When some critical value satisfies $\abs{P(\zeta)} < 1$, the lemniscate
components merge: two or more roots share a connected component.
The univalent branch $\psi_j$ can no longer be extended to all of $\D$,
as it encounters a branch point at $w = P(\zeta)$.

We present here our analysis of the degenerate case, including both
positive results and structural obstructions that clarify the difficulty of
the full EHP conjecture.

\begin{remark}\label{rem:squarefree-degenerate}
In Sections~\ref{sec:degenerate}--\ref{sec:numerics} we restrict attention
to squarefree polynomials (i.e., polynomials with simple roots) unless
explicitly stated otherwise.  This excludes repeated roots, for which the
formulas involving $P'(z_j)$, Blaschke zeros, and critical-point expansions
require separate treatment.
\end{remark}

%% --------------------------------------------------
\subsection{The conformal factorization}
%% --------------------------------------------------

Assume that $P$ is squarefree.
Let $z_{j^*}$ be a root satisfying
\[
  \abs{P'(z_{j^*})} \;=\; \min_{1 \leq j \leq n} \abs{P'(z_j)}.
\]
By Lemma~\ref{lem:fekete}, $\abs{P'(z_{j^*})} \leq n$.
Let $\mathcal{C}$ be the connected
component of $\Lambda_P(1)$ containing $z_{j^*}$.
Suppose $\mathcal{C}$ contains $k \geq 2$ roots (the degenerate case).
Let $\varphi \colon \D \to \mathcal{C}$ be the
Riemann map with $\varphi(0) = z_{j^*}$
(this exists since $\mathcal{C}$ is simply connected, as shown in
the proof of Lemma~\ref{lem:petals}).
The composition $B = P \circ \varphi \colon \D \to \D$ is a
degree-$k$ Blaschke product with $B(0) = 0$.

By the chain rule, $\varphi'(0) = B'(0)/P'(z_{j^*})$, and the
Koebe quarter theorem gives
\begin{equation}\label{eq:blaschke-inradius}
  \inrad(\mathcal{C}) \;\geq\; \frac{\abs{B'(0)}}{4\abs{P'(z_{j^*})}}.
\end{equation}
Since $B(w) = w \prod_{i=2}^k \frac{w - \beta_i}{1 - \bar\beta_i w}$
(the standard representation with $B(0) = 0$),
evaluating the derivative at $w = 0$ via the product rule gives
$\abs{B'(0)} = \prod_{i=2}^k \abs{\beta_i}$, where $\beta_2,
\ldots, \beta_k$ are the $\varphi$-preimages of the other roots in
$\mathcal{C}$.  Since $\abs{P'(z_{j^*})} \leq n$,
\[
  \inrad(\mathcal{C}) \;\geq\;
  \frac{\prod_{i=2}^k \abs{\beta_i}}{4n}.
\]
We emphasize that this bound applies to the component containing the
Fekete-selected root; for an arbitrary component, no such estimate is
available from the present arguments.

For the degenerate case, the Blaschke zeros $\beta_i$ may cluster
near the origin, causing $\prod\abs{\beta_i}$ to be small.
To capture the ``fatness'' of the merged component, one may instead
evaluate the conformal radius at a critical point of $B$.

%% --------------------------------------------------
\subsection{Conformal radius at critical points}
%% --------------------------------------------------

Let $w_c \in \D$ be a critical point of $B$ (i.e., $B'(w_c) = 0$),
and let $\zeta_c = \varphi(w_c)$ be the corresponding critical point
of $P$.  The Koebe bound at $w_c$ gives
\begin{equation}\label{eq:crit-koebe}
  \inrad(\mathcal{C}) \;\geq\;
  \frac{1}{4}\abs{\varphi'(w_c)}\,(1 - \abs{w_c}^2).
\end{equation}

To express $\abs{\varphi'(w_c)}$ in terms of the polynomial and Blaschke
data, we use a local expansion.

\begin{lemma}\label{lem:local-multiplicity}
Let $\varphi \colon \D \to \mathcal{C}$ be conformal and set
$B = P \circ \varphi$.  Let $w_c \in \D$ and $\zeta_c = \varphi(w_c)$.
Then $w_c$ is a critical point of $B$ of multiplicity $m \geq 1$
if and only if $\zeta_c$ is a critical point of $P$ of multiplicity $m$.
In that case,
$B^{(1)}(w_c) = \cdots = B^{(m)}(w_c) = 0$ and
\begin{equation}\label{eq:faa-di-bruno}
  B^{(m+1)}(w_c) \;=\; P^{(m+1)}(\zeta_c)\,\varphi'(w_c)^{m+1}.
\end{equation}
\end{lemma}

\begin{proof}
Write the local expansions
$\varphi(w) = \zeta_c + a(w - w_c) + O((w - w_c)^2)$ with
$a = \varphi'(w_c) \neq 0$, and
$P(z) = P(\zeta_c) + c\,(z - \zeta_c)^{m+1} + O((z - \zeta_c)^{m+2})$
where $c = P^{(m+1)}(\zeta_c)/(m+1)!$.
Substituting gives
$B(w) = P(\zeta_c) + c\,a^{m+1}(w - w_c)^{m+1} + O((w - w_c)^{m+2})$,
which is the claimed statement.
\end{proof}

\begin{corollary}\label{cor:simple-critical}
If $w_c$ is a simple critical point of $B$, then $\zeta_c = \varphi(w_c)$
is a simple critical point of $P$ and
\begin{equation}\label{eq:lhopital}
  \abs{\varphi'(w_c)}^2 \;=\;
  \frac{\abs{B''(w_c)}}{\abs{P''(\zeta_c)}}.
\end{equation}
\end{corollary}

\begin{proof}
Apply Lemma~\ref{lem:local-multiplicity} with $m = 1$.  Then
$B''(w_c) = P''(\zeta_c)\,\varphi'(w_c)^2$,
and taking absolute values gives~\eqref{eq:lhopital}.
\end{proof}

This cleanly separates the ``microscopic'' Blaschke geometry (the
numerator) from the ``macroscopic'' polynomial scaling (the denominator).

%% --------------------------------------------------
\subsection{A structural obstruction for $k \geq 2$}
\label{sec:obstruction}
%% --------------------------------------------------

For $k = 2$, the Blaschke product $B(w) = w \cdot \frac{w-\beta}{1-\bar\beta w}$
has a unique critical point $w_c$ inside $\D$.  Direct computation gives
\begin{equation}\label{eq:k2-identity}
  \abs{B''(w_c)}(1-\abs{w_c}^2)^2 \;=\; 2(1 - \abs{B(w_c)}^2).
\end{equation}
As the two roots in the component move further apart and the lemniscate
approaches a pinch point, the critical value satisfies $\abs{B(w_c)} \to 1$,
and consequently the hyperbolic second derivative degenerates to zero.
Thus the conformal radius estimate~\eqref{eq:crit-koebe} fails to provide
a universal lower bound even for two-root components.

For $k \geq 3$, the same obstruction persists.
Under the squarefree hypothesis (Remark~\ref{rem:squarefree-degenerate}),
the Blaschke product $B = P \circ \varphi$ has only simple zeros,
so the appropriate degree-$3$ family is
\[
  B_{a}(w) \;=\; w \cdot \frac{w - a}{1 - aw}
  \cdot \frac{w - 2a}{1 - 2aw}, \qquad a \in (0, \tfrac{1}{2}),
\]
which has three simple zeros at $0$, $a$, and $2a$.
For small $a$, the denominators are close to $1$ and
$B_a(w) \approx w(w-a)(w-2a)$, so
$B_a'(w) \approx 3w^2 - 6aw + 2a^2$
with two critical points at $w_c = a(1 \pm 1/\sqrt{3}) + O(a^2)$,
both of size $O(a)$.  A direct computation gives
\[
  \abs{B_a''(w_c)}(1-\abs{w_c}^2)^2 \;=\; 2\sqrt{3}\,a + O(a^2)
\]
at both critical points.

As $a \to 0$, the hyperbolic second derivative at \emph{both}
critical points tends to zero:
\[
  \max_{w_c}\, \abs{B_a''(w_c)}(1-\abs{w_c}^2)^2 \;\to\; 0.
\]
Consequently, the conformal radius estimate~\eqref{eq:crit-koebe}
degenerates, and the particular low-order Blaschke critical-point
invariants considered above do not by themselves yield a universal
bound for $k \geq 2$.

\begin{remark}
Geometrically, as $a \to 0$, the three zeros of $B_a$ coalesce at
the origin and the Blaschke product approaches $B_0(w) = w^3$.
In the limit, the critical point at $0$ has multiplicity $2$,
and Lemma~\ref{lem:local-multiplicity}
gives $\abs{\varphi'(0)}^3 = \abs{B'''(0)}/\abs{P'''(\zeta_c)} = 6/\abs{P'''(\zeta_c)}$,
which is bounded away from zero.  The obstruction is thus a
\emph{near-degeneracy} phenomenon: the Blaschke invariants degenerate
continuously as the zeros cluster, even though they
``reset'' at the exact moment of collision via the higher derivative.
This indicates that a proof for merged components cannot rely solely
on a uniform lower bound for one distinguished low-order Blaschke
derivative, though it does not rule out other conformal approaches.
\end{remark}

%% --------------------------------------------------
\subsection{An identity for $P''$ at critical points}
%% --------------------------------------------------

Despite the Blaschke obstruction, the polynomial structure provides
useful identities.  Differentiating the logarithmic derivative
$P'(z)/P(z) = \sum_{j=1}^n 1/(z-z_j)$ and evaluating at a critical
point $\zeta$ where $P'(\zeta) = 0$ and $P(\zeta) \neq 0$
(i.e., $\zeta$ is not a repeated root), we obtain
\begin{equation}\label{eq:Ppp-identity}
  P''(\zeta) \;=\; -P(\zeta) \sum_{j=1}^n \frac{1}{(\zeta - z_j)^2}.
\end{equation}

This identity cleanly separates local and global contributions:
roots near $\zeta$ dominate the sum via their large $1/(\zeta-z_j)^2$
terms, while distant roots contribute negligibly.

\begin{lemma}[Dual Fekete bound]\label{lem:dual-fekete}
Let $P$ be a monic polynomial of degree $n$ with all roots in
$\overline{\D}$.  If $P'$ has a repeated zero, then
$\min_{\zeta:\,P'(\zeta)=0}\abs{P''(\zeta)} = 0$.
Otherwise, if $\zeta_1, \ldots, \zeta_{n-1}$ are the distinct
critical points of $P$, then
\begin{equation}\label{eq:dual-fekete}
  \min_{1 \leq i \leq n-1} \abs{P''(\zeta_i)} \;\leq\; n(n-1).
\end{equation}
\end{lemma}

\begin{proof}
If $P'$ has a repeated zero, the first claim is immediate.
Assume now that $P'$ has distinct zeros $\zeta_1, \ldots, \zeta_{n-1}$.
Since $P'(z) = n\prod_{j=1}^{n-1}(z - \zeta_j)$, we have
$P''(\zeta_i) = n\prod_{j \neq i}(\zeta_i - \zeta_j)$.
By the Gauss--Lucas theorem, each $\zeta_i$ lies in $\overline{\D}$.
Applying Lemma~\ref{lem:fekete} to the $n-1$ points
$\zeta_1, \ldots, \zeta_{n-1}$ gives
$\min_i \prod_{j \neq i}\abs{\zeta_i - \zeta_j} \leq n-1$,
and therefore $\min_i \abs{P''(\zeta_i)} \leq n(n-1)$.
\end{proof}

This ``dual Fekete bound'' establishes that the critical points are
subject to the same macroscopic energy constraints as the roots.

%% --------------------------------------------------
\subsection{Numerical evidence}
\label{sec:numerics}
%% --------------------------------------------------

We computed the inradius for
$P \in \Pn(\overline{\D})$ with $n = 10$ across several
hundred random configurations: roots uniformly distributed on $\T$,
roots in the interior of $\D$, clustered configurations, and
adversarial gradient searches.  (The inradius was approximated by
evaluating $\abs{P}$ on a fine polar grid and performing a binary search
on the radius of the largest inscribed disc.)
In all cases tested, the minimum value of
$n \cdot \inrad(\Lambda_P)$ was attained by roots of unity,
with $n \cdot \inrad(\Lambda_{z^{10}-1}) \approx 1.26$.

Degenerate configurations (with roots in the interior) consistently
exhibited \emph{larger} inradii---typically $n \cdot \inrad \geq 4$.
Even among non-degenerate configurations on $\T$, the roots of unity
appeared extremal; perturbing any root away from the uniform spacing
increased the inradius.

These observations are consistent with the hypothesis that the inradius
minimizer is non-degenerate (and in fact equals $z^n - 1$), which would
immediately imply the full EHP conjecture via Theorem~\ref{thm:main}.
Since these computations play no role in the proofs, we record
them only as motivation.

%% ====================================================================
\section{Comparison with KLR approach}
%% ====================================================================

The KLR proof of \eqref{eq:KLR} proceeds by:
\begin{enumerate}[(1)]
  \item Using the Nazarov--Polterovich--Sodin theorem to show
    $\operatorname{Area}(\Lambda_P) \geq c/\log n$;
  \item Bounding the perimeter via a Fryntov--Nazarov formula:
    $L(\Lambda_P) \leq 4n\sqrt{\pi \cdot \operatorname{Area}(\Lambda_P)}$;
  \item Applying an isoperimetric inequality:
    $\operatorname{Area} \leq 18\pi \cdot \inrad \cdot L$.
\end{enumerate}
Combining these gives $\inrad \geq c\sqrt{\operatorname{Area}}/n \geq c/(n\sqrt{\log n})$.
The $\sqrt{\log n}$ loss is structural: it arises from the square root
in the area-to-inradius conversion via the perimeter bound.

Our approach is fundamentally different: we never estimate area or
perimeter.  Instead, we work directly with the \emph{conformal structure}
of the polynomial map $P \colon \C \to \C$, using the branched covering
structure to construct large inscribed discs.  This explains why we obtain
the sharp $1/n$ order without logarithmic losses---the Koebe theorem
provides a direct pointwise estimate, bypassing the global integral
estimates that introduce the square root.

%% ====================================================================
\section{Reduction of the full conjecture}
%% ====================================================================

We summarize the implications of our analysis for the full EHP conjecture.

\begin{lemma}\label{lem:minimizer}
For each $n \geq 1$, there exists $P^* \in \Pn(\overline{\D})$ such that
$\inrad(\Lambda_{P^*}(1)) = \rho_n$.
\end{lemma}

\begin{proof}
The parameter space $\overline{\D}^n$ (identifying a monic polynomial
with its $n$ roots) is compact.  It suffices to show that
$a \mapsto \inrad(\Lambda_{P_a}(1))$ is continuous.
All lemniscates are contained in $\{z : \abs{z} \leq 2\}$
(since $\abs{z} \geq 2$ and $\abs{a_j} \leq 1$ imply
$\abs{P_a(z)} \geq 1$), so the inradius is bounded.
Lower semicontinuity follows from uniform convergence of polynomials
on compact sets: if $B(z_0,r) \subset \Lambda_{P_a}(1)$, then
$\max_{\overline{B(z_0,r)}}\abs{P_a} < 1$, and for $a'$ sufficiently
close to $a$, $\abs{P_{a'}} < 1$ on $\overline{B(z_0,r)}$ as well.
Upper semicontinuity follows similarly: if $\abs{P_{a^{(m)}}} < 1$
on $B(z_m, r_m)$ with $z_m \to z$ and $r_m \to L$, then
$\abs{P_a} \leq 1$ on $B(z,r)$ for any $r < L$, hence
$\abs{P_a} < 1$ by the maximum principle (since $P_a$ is nonconstant).
A continuous function on a compact set attains its infimum.
\end{proof}

\begin{proposition}\label{prop:reduction}
The following implications hold:
\begin{enumerate}[(i)]
  \item If some inradius minimizer in $\Pn(\overline{\D})$
    is non-degenerate, then the EHP conjecture holds:
    $\rho_n \geq 1/(4n)$.
  \item For every $P \in \Pn(\overline{\D})$, if there exists a
    connected component $\mathcal{C}$ of $\Lambda_P(1)$ containing
    $k$ roots with $\inrad(\mathcal{C}) \geq c/n$, then
    $\inrad(\Lambda_P(1)) \geq c/n$.
\end{enumerate}
In particular, one sufficient route to the full EHP conjecture is to
prove that some inradius minimizer is non-degenerate.
\end{proposition}

\begin{proof}
For~(i): if the minimizer $P^*$ is non-degenerate, then
Theorem~\ref{thm:main} gives
$\rho_n = \inrad(\Lambda_{P^*}) \geq 1/(4n)$.
Statement~(ii) is immediate since
$\inrad(\Lambda_P) \geq \inrad(\mathcal{C})$.
\end{proof}

\begin{remark}
The converse of~(i) does not hold a priori: the minimizer could
theoretically be degenerate while still satisfying $\inrad \geq c/n$.
However, no proof of non-degeneracy for the minimizer is currently known,
and this remains a conjecture supported by the numerical evidence of
Section~\ref{sec:numerics}.
\end{remark}

Our analysis identifies two structural barriers to proving the
full conjecture directly for degenerate components:
\begin{enumerate}[(i)]
  \item The examples of Section~\ref{sec:obstruction} show that
    the particular low-order Blaschke critical-point invariants
    considered here can degenerate for
    $k \geq 2$, so this line of argument does not by itself yield
    a universal bound.
  \item The inradius is not monotone under root perturbation
    (unlike area, which decreases under radial pushing by the
    zero-pushing lemma of \cite{KLR25}), preventing variational
    proofs of non-degeneracy for the minimizer.
\end{enumerate}

We propose three possible routes to the full conjecture:
\begin{enumerate}
  \item \textbf{Non-degeneracy of minimizers.}
    Prove that the inradius minimizer is non-degenerate by
    non-variational methods (e.g., direct analysis of the
    extremal conditions).  The numerical evidence of
    Section~\ref{sec:numerics} is consistent with this.

  \item \textbf{Curvature-based width bounds.}
    The total curvature of a degree-$n$ lemniscate boundary
    satisfies $\int \abs{\kappa}\,ds \leq 2\pi n$.
    This prevents arbitrarily thin, wiggly components and may
    yield direct width bounds that bypass the area-perimeter chain.

  \item \textbf{Capacity-diameter-inradius inequalities.}
    For a connected monic lemniscate $\{z : \abs{P(z)} < 1\}$
    of degree $n$ with $\Cap = 1$ and all roots in $\overline{\D}$,
    prove $\inrad \geq c/n$ using potential-theoretic methods.
    The capacity constraint is the key rigidity that separates
    polynomial lemniscates from general harmonic sublevel sets
    (for which the analogous statement is false).
\end{enumerate}

%% ====================================================================
%% References
%% ====================================================================

\begin{thebibliography}{99}

\bibitem{Duren83}
P.~L.~Duren, \emph{Univalent Functions},
Grundlehren der mathematischen Wissenschaften \textbf{259},
Springer-Verlag, New York, 1983.

\bibitem{EHP58}
P.~Erd\H{o}s, F.~Herzog, and G.~Piranian,
``Metric properties of polynomials,''
\emph{J.\ Analyse Math.}\ \textbf{6} (1958), 125--148.

\bibitem{KLR25}
M.~Krishnapur, E.~Lundberg, and K.~Ramachandran,
``On the area of polynomial lemniscates,''
arXiv:2503.18270 [math.CV], March 2025.

\bibitem{Pom61}
Ch.~Pommerenke,
``On metric properties of complex polynomials,''
\emph{Michigan Math.\ J.}\ \textbf{8} (1961), 97--115.

\bibitem{Tao25}
T.~Tao,
``The maximal length of the Erd\H{o}s--Herzog--Piranian lemniscate in high degree,''
arXiv:2512.12455, December 2025.

\end{thebibliography}

\end{document}
